\documentclass[10pt,a4paper]{amsart}
\usepackage[margin=19mm]{geometry}
\usepackage{amsmath,amssymb,mathtools}
\usepackage[scr=rsfs,cal=euler]{mathalfa}
\usepackage{xfrac}
\usepackage[unicode,hidelinks,bookmarks=false]{hyperref}
\newcommand{\ie}{i.e.\ }
\newcommand{\cf}{cf.\ }
\newcommand{\resp}{respectively\ }
\newcommand{\Subsection}{\subsection}
\providecommand{\nobreakdash}{\nobreak\hbox{-}}
\setlength{\emergencystretch}{2em}
\multlinegap0pt
\usepackage{bm}
\usepackage{bbm}
\usepackage{dsfont}
\usepackage{xspace}
\usepackage{cancel}
\usepackage{mathtools}
\usepackage{amsthm}
\makeatletter
\@ifundefined{theorem}{\newtheorem{theorem}{Theorem}}{}
\@ifundefined{lemma}{\newtheorem{lemma}[theorem]{Lemma}}{}
\makeatother
\newcommand{\QQ}{\mathbb{Q}}
\newcommand{\ZZ}{\mathbb{Z}}
\newcommand{\lrp}[1]{\!\left(#1\right)}
\renewcommand{\Im}{\mathrm{Im}}
\title[Only finitely many modular Jacobians are supersingular modul]{Only finitely many modular Jacobians are supersingular modulo a given prime}
\author{Aarya J. Kumar, Sargam Mondal, Erick Ross, Hui Xue}
\date{Source: arXiv:2609.17819v1 (CC BY 4.0)}
\begin{document}
\maketitle

\noindent\emph{Source and licence.} Only finitely many modular Jacobians are supersingular modulo a given prime. Version arXiv:2609.17819v1 (CC BY 4.0), distributed under the Creative Commons Attribution 4.0 International licence, as recorded in the arXiv licence metadata for this exact version and shown on the abstract page https://arxiv.org/abs/2609.17819v1. Full rights notice: licenses/arxiv-2609.17819-CC-BY-4.0.txt.\par\smallskip

\noindent\textit{Editorial scope.} Counted unit: the Lemma whose source label is \texttt{lem:conjugates-shape} in the article (source0.tex lines 563--567, source label \texttt{lem:conjugates-shape}), delivered together with the article's own complete original proof of it (source0.tex lines 569--685, one \texttt{proof} environment of 117 lines). No mathematics is written, repaired or re-derived: statement and proof are the article's own text line for line. Required context retained uncounted: none. Every definition, display and label of those lines is retained unchanged. Omitted from this package and not redistributed: 1--562, 568--568, 686--1567. Nothing inside a retained excerpt is omitted or altered: every retained body is delivered exactly as the article's lines are written, so no omission receipt of any kind accompanies this package. Citations: the retained excerpts carry no citation command at all, so no bibliography entry of the article is transcribed and no reference is redistributed. Reference adaptation: none was needed and none was made; every reference inside the delivered unit resolves inside it, so no unresolved reference or citation is produced. Printed slips kept literal: no printed word, symbol, hypothesis or number of the counted statement or of its proof was altered. Layout and class: the article's own class is \texttt{amsart} with packages including fullpage, amsmath, amsthm, amssymb, amscd, enumerate, enumitem, array, float, bbm, bm, stmaryrd, comment, mathtools; the wrapper is the lane's pinned \texttt{amsart} editorial wrapper at 10pt on A4, which restates the theorem-like environments the retained text needs and adds the article's own macro definitions that the retained text needs (\texttt{\textbackslash Im}, \texttt{\textbackslash QQ}, \texttt{\textbackslash ZZ}, \texttt{\textbackslash lrp}), with extra wrapper packages (bm, bbm, dsfont, xspace, cancel, mathtools). The delivered numbering is the unit's own. Assets: no retained span contains a figure environment, a graphics-inclusion command, a PDF-page inclusion, a table or a TikZ picture; no image file is copied into this package, the package declares no asset dependency and no image file is redistributed with it. Licence: version arXiv:2609.17819v1 of this article is distributed under the Creative Commons Attribution 4.0 International licence, as recorded in the arXiv licence metadata for this exact version and shown on the abstract page https://arxiv.org/abs/2609.17819v1. A full rights notice is delivered at licenses/arxiv-2609.17819-CC-BY-4.0.txt. Characters: every retained excerpt is pure ASCII, so the delivered engine, XeLaTeX with its default Latin Modern text font, reports no missing character.
\begin{lemma}\label{lem:conjugates-shape}
For every prime $p$ and every
primitive $n$-th root of unity $\zeta$, 
$\frac{1}{\sqrt p} \operatorname{Conj}_{\operatorname{Gal}(\overline{\QQ}/\QQ)}(\sqrt p\,\zeta)$ is equal to a finite disjoint union of $\Phi_n^+$-imitators.
\end{lemma}

\begin{proof}
Let $L=\QQ(\sqrt p,\zeta)$. Since every Galois conjugate of $\sqrt p\,\zeta$ lies in $L$, first observe that
\begin{align}
    \operatorname{Conj}_{\operatorname{Gal}(\overline{\QQ}/\QQ)}(\sqrt p\,\zeta)=\big\{\tau(\sqrt p\,\zeta):\tau\in\operatorname{Gal}(L/\QQ)\big\} \qquad \text{(at least, as sets)}.
\end{align} 
Then we break into cases.

{\bf Case 1: $\sqrt p \notin \QQ(\zeta)$}. \\ 
Since $\sqrt p\notin\QQ(\zeta)$, we have that $\QQ(\sqrt p)\cap\QQ(\zeta)=\QQ$, so 
\[
\operatorname{Gal}(L/\QQ)\cong\operatorname{Gal}(\QQ(\sqrt p)/\QQ)\times\operatorname{Gal}(\QQ(\zeta)/\QQ).
\] 
This means that the automorphisms in $\mathrm{Gal}(L/\QQ)$ are precisely those of the form $(\rho,\tau_a)$, where
$\rho(\sqrt p)=\pm\sqrt p$ and 
$\tau_a(\zeta)=\zeta^a$ for $a\in(\ZZ/n\ZZ)^\times$. 
Hence
\begin{align}
    \frac{1}{\sqrt p} \operatorname{Conj}_{\operatorname{Gal}(\overline{\QQ}/\QQ)}(\sqrt p\,\zeta)
    &= \frac{1}{\sqrt p} \Big\{\tau(\sqrt p\,\zeta):\tau\in\operatorname{Gal}(L/\QQ)\Big\} \\
    &= \Big\{ \pm \zeta^a : a\in(\ZZ/n\ZZ)^\times \Big\} \\
    &= \Phi_n \cup -\Phi_n \\
    &= 
    \begin{cases}
        \Phi_n \sqcup -\Phi_n & \text{if } 4 \nmid n \\
        \Phi_n & \text{if } 4 \mid n.
    \end{cases}
\end{align}

This implies the desired result since $\Phi_n$ is a finite disjoint union of $\Phi_n^+$-imitators; namely
\begin{align}
    \Phi_n = 
    \begin{cases}
        \{ z : z \in \Phi_n^+\} 
        & \text{if } n=1,2 \\
        \{z : z \in\Phi_n^+\}\sqcup\{z^{-1} : z \in \Phi_n^+\}
        & \text{if } n \ge 3.
    \end{cases} 
\end{align}


\textbf{Case 2: $\sqrt p \in \QQ(\zeta)$.}\\ 
In this case, we have
\[
\operatorname{Gal}(L/\QQ)
=
\operatorname{Gal}(\QQ(\zeta)/\QQ),
\] 
and so the automorphisms in $\mathrm{Gal}(L/\QQ)$ are precisely those of the form $\tau_a$, where
$\tau_a(\zeta)=\zeta^a$ for $a\in(\ZZ/n\ZZ)^\times$. Now, each of these automorphisms must send $\sqrt p$ to $\pm\sqrt p$. This then implies that 
$\tau_a(\sqrt p)=\chi(a)\sqrt p$, where $\chi$ is some quadratic Dirichlet character modulo $n$.


Consequently,
\begin{align}
    &\frac{1}{\sqrt p} \operatorname{Conj}_{\operatorname{Gal}(\overline{\QQ}/\QQ)}(\sqrt p\,\zeta) \\
    &=
    \frac{1}{\sqrt p} \Big\{\tau(\sqrt p\,\zeta):\tau\in\operatorname{Gal}(L/\QQ)\Big\} \\
    &=
    \frac{1}{\sqrt p} \Big\{\tau_a(\sqrt p) \tau_a(\zeta) : a \in (\ZZ/n\ZZ)^{\times}\Big\} \\
    &=
    \Big\{\chi(a)\,\zeta^a : a \in (\ZZ/n\ZZ)^{\times}\Big\} \\
    &=
    \Big\{\chi(a)\,\zeta^a, \chi(n-a) \zeta^{n-a} : a \in (\ZZ/n\ZZ)^{\times},\ 0 < a < \tfrac{n}{2} \Big\} \\
    &=
    \Big\{\chi(a)\,\zeta^a : a \in (\ZZ/n\ZZ)^{\times},\ 0 < a < \tfrac{n}{2} \Big\} 
    \cup 
    \Big\{\chi(-a) \zeta^{-a} : a \in (\ZZ/n\ZZ)^{\times},\ 0 < a < \tfrac{n}{2} \Big\}
    \\ 
    &= 
    \begin{cases}
        \Big\{\chi(a)\,\zeta^a : a \in (\ZZ/n\ZZ)^{\times},\ 0 < a < \tfrac{n}{2} \Big\}
        &\hspace{-45mm}
        \quad\text{if } 4\mid n,\ \chi(1-\tfrac{n}{2})=-1,  \\
        \Big\{\chi(a)\,\zeta^a : a \in (\ZZ/n\ZZ)^{\times},\ 0 < a < \tfrac{n}{2} \Big\} 
        \sqcup
        \Big\{\chi(-a)\,\zeta^{-a} : a \in (\ZZ/n\ZZ)^{\times},\ 0 < a < \tfrac{n}{2} \Big\}
        \label{eqn:temp-conj-formula-disjoint-union}
        \\
        &\hspace{-17.5mm} \text{otherwise.}  
    \end{cases}  
\end{align}
Here, the last step follows from the equivalence
\begin{align} \label{eqn:temp-twisted-roots-of-unity-overlap-equiv-condition}
    &\quad\ \ \,  \chi(a) \zeta^a = \chi(-a') \zeta^{-a'}  \\
    &\text{iff}\quad
    \chi(-aa') =\zeta^{a+a'}  
    & \text{(since $\chi$ is quadratic)} \\
    &\text{iff}\quad
    \chi(-aa') = -1
    \text{ and }
    a+a' = \tfrac{n}{2} 
    &\text{(since $\chi$ is quadratic and $0 < a,a' < \tfrac{n}{2}$)}
    \\
    &\text{iff}\quad
    4 \mid n,\ \chi\lrp{-a(\tfrac{n}{2} - a)}=-1, \text{ and } 
    a' = \tfrac{n}{2}-a 
    &\text{(since $\tfrac{n}{2} - a \in (\ZZ/n\ZZ)^\times$ only when $4 \mid n$)} 
    \\
    &\text{iff}\quad
    4 \mid n,\ \chi\lrp{1-\tfrac{n}{2}}=-1, \text{ and } 
    a' = \tfrac{n}{2}-a. \\
    & &
    \hspace{-80mm}
    \Big(\text{since $\chi(-a(\tfrac{n}{2}-a)) = \chi(a^2 - a^2 \tfrac{n}{2} + a^2 \tfrac{n}{2} - a \tfrac{n}{2}) = \chi(a^2 (1-\tfrac{n}{2}) + \tfrac{a(a-1)}{2}n) = \chi(1 - \tfrac{n}{2})$}\Big). 
\end{align}
Also, note that all the elements are distinct within each set of \eqref{eqn:temp-conj-formula-disjoint-union}, since $\zeta^{a-a'} \ne \pm 1$ for $a \ne a'$. 

Finally, observe that $\big\{\zeta^a : a \in (\ZZ/n\ZZ)^{\times},\ 0 < a < \tfrac{n}{2} \big\}$ is a $\Phi_n^+$-imitator since
\begin{align}
    \Big\{ \big(\zeta^a\big)^{\mathrm{sgn}\,\Im(\zeta^a)} :  a \in (\ZZ/n\ZZ)^{\times},\ 0 < a < \tfrac{n}{2} \Big\} = \Phi_n^+.\label{eqn:temp-set-equal-to-Phi_n+} 
\end{align}
Hence $\big\{\chi(a)\,\zeta^a : a \in (\ZZ/n\ZZ)^{\times},\ 0 < a < \tfrac{n}{2} \big\}$ 
and 
$\big\{\chi(-a)\,\zeta^{-a} : a \in (\ZZ/n\ZZ)^{\times},\ 0 < a < \tfrac{n}{2} \big\}$ 
are also $\Phi_n^+$-imitators, so \eqref{eqn:temp-conj-formula-disjoint-union}
yields the desired result.
\end{proof}

\end{document}
